\section{Prácticas de despliegue y monitoreo}
\subsection{Integración con CI/CD}
Graham sugiere combinar el Agent con SWE-bench dentro del pipeline de CI: ejecutar el conjunto Lite todos los días o en cada commit, y generar un issue con las muestras que fallan para que el equipo de modelos las estudie. Los scripts automatizados pueden ejecutar Lite/Full en paralelo en varios entornos, y subir los registros para que las personas los revisen.

\begin{lstlisting}[language=bash]
#!/usr/bin/env bash
set -ex
cd /workspace/swe-bench
python run_swe_bench.py --split lite --model codegen
python run_swe_bench.py --split full --tests smoke
tar czf swe-bench-results.tar.gz outputs/
gzip -c outputs/logs/latest.log > swe-bench-results/logs/latest.log.gz
notify-team --subject "SWE-bench run finished" --files swe-bench-results.tar.gz
\end{lstlisting}

El \texttt{notify-team} del script puede conectarse a Slack o correo electrónico, y enviar junto el diff de la falla y las entradas/salidas generales.

\begin{knowledgebox}{Estrategia de ejecución de CI/CD}
Divide la ejecución del Agent en «precalentamiento, ejecución y reproducción»: en la etapa de precalentamiento se reconstruyen el repositorio y las dependencias, en la etapa de ejecución se ejecuta el Agent, y en la etapa de reproducción se suben los registros y el diff para que los desarrolladores los revisen.
\end{knowledgebox}

\subsection{Experimentos a pequeña escala}
Antes de ampliar el despliegue del Agent, primero realiza 1 o 2 tareas piloto en el repositorio \texttt{sandbox}: usa el Agent para generar el patch, pero revísalo solo en la red interna, y mantén cada ronda de ejecución en menos de 15 minutos para asegurar que sea controlable.

\begin{importantbox}{Recomendaciones para experimentos de corto plazo}
Conserva puntos de verificación humanos: permite que el Agent modifique menos de 3 archivos, y toda operación debe generar un patch en \texttt{git} que una persona confirme después.
\end{importantbox}

\begin{warningbox}{Evitar la contaminación de la rama principal}
No dejes que el Agent haga push directamente a `main`/`master`; debe generar el patch en una rama feature y fusionarla hasta que termine la revisión.
\end{warningbox}

\subsection{Resumen de la sección}
\begin{itemize}
    \item Integrar el Agent en el CI permite detectar regressions de manera periódica, y retroalimentar los trace de las fallas.
    \item Los experimentos a pequeña escala en la red interna pueden reducir el riesgo de iteración y acumular experiencia.
\end{itemize}

